\section{Introduction}
\label{sec:introduction}

Magic squares are among the oldest objects of recreational mathematics: the $3\times3$ Lo Shu appears in Chinese sources more than two millennia old, a $4\times4$ square is engraved in D\"urer's \emph{Melencolia~I} (1514), and systematic constructions for every order have been known since the seventeenth century \cite{andrews1960magic,pickover2002zen}.
The modern literature treats them from several directions: enumeration and the geometry of the associated polytopes \cite{stanley1973magic,ehrhart1973magic,ahmed2003polyhedral,beck2006insideout,beck2010enumeration}, spectral and matrix properties \cite{loly2009spectra,nordgren2012properties,nordgren2014spectra,cameron2013signatura}, and, in work of Loly and collaborators, physical invariants such as the moment of inertia and the multipole moments \cite{loly2004invariance,rogers2004inertia,rogers2005multipole}.

\begin{definition}[Magic square]
\label{def:magic_square}
An \emph{arrangement} of order $n$ is a bijection $S=(s_{ij})$ from the cells of the $n\times n$ grid onto $\{1,\dots,n^2\}$.
It is a \emph{magic square} if every row, every column and both main diagonals sum to the magic constant $M(n)=n(n^2+1)/2$, and \emph{semi-magic} if only the rows and columns do.
\end{definition}

This paper studies a three-dimensional representation of an arrangement: place the centred entry $s_{ij}-(n^2+1)/2$ as a height over the centred cell $(j-\frac{n-1}2,\ \frac{n-1}2-i)$ and take the convex hull of the $n^2$ points, the \emph{Magic Gem} of $S$ (\Cref{sec:gem}).
The representation is invariant under the dihedral group $D_4$ of the square, so each essentially different magic square has a single gem (\Cref{prop:symmetry_invariance}); an interactive web application displays the $880$ essentially different order-four squares and renders the gem of a selected one (points, convex hull, vectors from the centre, coordinate axes; rotatable, with orthographic views) at \url{https://kylemath.github.io/MagicGemWebpage/}.
The convex hull itself is the object of Sections~\ref{sec:gem-geometry-880}--\ref{sec:siamese-gems} (the geometry of the $880$ order-four gems and of the Siamese gems); everywhere else the object is the height function $z$.
Once the entries are read as a function on the grid, every line-sum condition ``row $k$ sums to $M(n)$'' becomes ``$z$ is orthogonal to the centred indicator of row $k$'', i.e.\ $\Cov(R_k,Z)=0$ for the empirical covariance over the $n^2$ cells.
A magic square is therefore an arrangement whose height function is orthogonal to a fixed $2n$-dimensional space $W$ of test functions, and the natural \emph{energy} $\Efull(S)=\sum_L\Cov(\mathbf 1_L,Z)^2$ over the $2n+2$ lines vanishes exactly on magic squares (\Cref{lem:zero_set}).
We are explicit that this is a reformulation of the definition and not a theorem about magic squares (\Cref{rem:tautology}); its value is that it suggests questions whose answers are not visible in the line-sum formulation, and the paper is organised around those answers.

\paragraph{Contributions.}
\begin{enumerate}\setlength\itemsep{1pt}
\item \emph{The moment hierarchy} (\Cref{prop:hierarchy}, \Cref{thm:perm_threshold}). Replacing the line indicators by the polynomial test functions $x^m,y^m$ gives, for each degree $k$, the conditions $\Cov(X^m,Z)=\Cov(Y^m,Z)=0$ for $m\le k$. For real-valued height functions these are equivalent to equal row and column sums if and only if $k\ge n-1$; this is the invertibility of a Vandermonde matrix and we state it as a proposition. The result of the section is \Cref{thm:perm_threshold}: for permutations of $1,\dots,n^2$, integrality and the bound on line sums force the threshold $k^*(n)$ down to at most $n-2$ for every $n\ge14$, whereas $k^*(n)=n-1$ for $3\le n\le13$ (computer-verified witnesses, \Cref{rem:kstar_search}). The ``low-mode'' energy $\Elow$ of earlier versions of this paper is the degree-one truncation. Loly's invariance of the moment of inertia \cite{loly2004invariance}, which Loly derived from the row and column sums alone and so stated for semi-magic squares, is the single condition $\Cov(Z,X^2+Y^2)=0$: implied by the degree-two conditions, but much weaker than them (\Cref{cor:loly}). The same identity gives the null distribution of the moment of inertia of a random arrangement (\Cref{cor:loly_null}).
\item \emph{Exact statistics of a random arrangement} (\Cref{prop:random_energy_moments}). $\E[\Elow]=(n^2+1)(n^4-n^2+12n-12)/(72n^2)$ and $\E[\Efull]=(n^4-1)/(6n^2)$, with exact variances; $\Elow/\E[\Elow]\to\mathrm{Exp}(1)$ by Hoeffding's combinatorial central limit theorem. Monte Carlo with $10^8$ samples per order agrees for $n\le 8$ to within $5\times10^{-4}$ relative error.
\item \emph{The single-swap formula} (\Cref{thm:swap}). Exchanging two entries of a magic square that differ by $d$ raises $\Efull$ by exactly $d^2m/n^4$, where $m\in\{2,\dots,6\}$ counts the lines separating the two cells; every magic square is a strict local minimum, with minimum gap between $2/n^4$ and $6/n^4$.
\item \emph{The low-mode class for $n=4$} (\Cref{prop:L4}). Exactly $14{,}639{,}440$ of the $16!$ arrangements satisfy $\Elow=0$, $7{,}040$ of them magic; the count is a meet-in-the-middle enumeration and agrees with a Gaussian local-limit estimate to within $20\%$.
\item \emph{Geometry of the $880$ order-four gems} (\Cref{sec:gem-geometry-880}). $8$ to $16$ hull vertices, $39$ $f$-vectors, $352$ combinatorial types and $616$ congruence classes, which are exactly the orbits of the complement map $S\mapsto17-S$; no two gems have the same hull-vertex set. The Siamese squares of every odd order $n\ge5$ have gems with exactly twelve vertices (\Cref{prop:siamese-12}, computer-assisted proof).
\end{enumerate}

\paragraph{What is and is not new.}
The polyhedral geometry of magic squares is classical: nonnegative integer matrices with all line sums equal to $t$ are the lattice points of the $t$-th dilate of the Birkhoff polytope and their number is a polynomial in $t$ \cite{stanley1973magic,ehrhart1973magic}; the count of magic squares with distinct entries is an Ehrhart quasi-polynomial of an inside-out polytope \cite{beck2006insideout,beck2010enumeration}; the cones of magic squares and cubes have Hilbert-basis descriptions \cite{ahmed2003polyhedral}.
The Magic Gem is a different object: not the polytope whose lattice points are the squares, but a polyhedron attached to one square.
The orthogonality reading of the constraints is elementary and implicit in Loly's work, and the hierarchy itself is elementary: that the axial moments $\sum_{ij}s_{ij}x_j^m$ of a semi-magic object vanish for every $m$ is, in physical language, the vanishing of its multipole moments, as Rogers and Loly observed for the dipole and quadrupole moments of a magic cube \cite{rogers2005multipole}.
What is new is the converse threshold and its permutation refinement (\Cref{thm:perm_threshold}), the exact random-arrangement moments, the swap formula, the exact low-mode count, and the exact geometry of the $880$ gems.
We make no claim that the energy reformulation by itself is a result.

\paragraph{Corrections relative to version 4 (arXiv:2512.09170v4).}
\begin{itemize}\setlength\itemsep{0pt}
\item[(a)] The complete energy used there omitted one row and one column and was not $D_4$-invariant; all statistics now use the symmetric sum \eqref{eq:Efull}.
\item[(b)] The sampled statistics (means $4.44$, $12.04$, $26.00$), the ``landscape entropy'' and the quadratic fits had the diagonal terms weighted by $n^2$; withdrawn, replaced by \Cref{prop:random_energy_moments}.
\item[(c)] The ``trimodal'' swap gaps and the ``$1/n^2$'' minimum-gap scaling were artifacts: gaps are $d^2m/n^4$ and $\Delta_{\min}=\Theta(n^{-4})$ with constants $2$, $6$ (\Cref{cor:local_min}).
\item[(d)] The $24$ local minima of $\Efull$ for $n=3$ referred to the asymmetric energy; the symmetric one has $16$ (\Cref{sec:n3_exhaustive}).
\item[(e)] ``$\Elow=0$ characterises magic (or semi-magic) squares'' is withdrawn (\Cref{rem:four_term_insufficient}, \Cref{prop:L4}); $P(\Elow=0)$ for $n=4$ is $7.0\times10^{-7}$, not ``below $10^{-8}$'', and the statement that $6\times10^7$ samples for $n=4$ and $4\times10^8$ for $n=5$ contained no zero is unsupported by the archived data, which hold $10^5$ samples per order; the cause cannot be reconstructed from the repository (\Cref{sec:lowmode_n4}).
\item[(f)] Table~2 of v4 (``typical'' $16$ hull vertices for $n=5$, volume $95.0$ for $n=4$) is replaced by exact values for named squares (\Cref{sec:gem-geometry-880}); the hardware statement is corrected (\Cref{sec:appendix}).
\end{itemize}

\Cref{sec:background} collects the facts we use; \Cref{sec:gem} defines the gem; \Cref{sec:theory} contains the projection lemma, the hierarchy and the swap formula; \Cref{sec:results} the computations; \Cref{sec:discussion} what the viewpoint does and does not buy, local search, and open questions; the appendices the methods, the variance closed forms, the witnesses and the remaining steps of the Siamese proof.
